\documentclass[10pt,a4paper]{amsart}
\usepackage{amsmath,amssymb,amsthm,graphicx}
\usepackage{verbatim}
\newenvironment{reference}{\comment}{\endcomment}
%\newenvironment{reference}{}{}
\newenvironment{slogan}{\comment}{\endcomment}
\newenvironment{history}{\comment}{\endcomment}
% For commutative diagrams we use Xy-pic
\usepackage[all]{xy}

% We use 2cell for 2-commutative diagrams.
\xyoption{2cell}
% Theorem environments.
%
\theoremstyle{plain}
\newtheorem{theorem}[subsection]{Theorem}
\newtheorem{proposition}[subsection]{Proposition}
\newtheorem{lemma}[subsection]{Lemma}

\theoremstyle{definition}
\newtheorem{definition}[subsection]{Definition}
\newtheorem{example}[subsection]{Example}
\newtheorem{exercise}[subsection]{Exercise}
\newtheorem{situation}[subsection]{Situation}

\theoremstyle{remark}
\newtheorem{remark}[subsection]{Remark}
\newtheorem{remarks}[subsection]{Remarks}

\numberwithin{equation}{subsection}

% Macros
%
\def\lim{\mathop{\mathrm{lim}}\nolimits}
\def\colim{\mathop{\mathrm{colim}}\nolimits}
\def\Spec{\mathop{\mathrm{Spec}}}
\def\Hom{\mathop{\mathrm{Hom}}\nolimits}
\def\Ext{\mathop{\mathrm{Ext}}\nolimits}
\def\SheafHom{\mathop{\mathcal{H}\!\mathit{om}}\nolimits}
\def\SheafExt{\mathop{\mathcal{E}\!\mathit{xt}}\nolimits}
\def\Sch{\mathit{Sch}}
\def\Mor{\mathop{\mathrm{Mor}}\nolimits}
\def\Ob{\mathop{\mathrm{Ob}}\nolimits}
\def\Sh{\mathop{\mathit{Sh}}\nolimits}
\def\NL{\mathop{N\!L}\nolimits}
\def\CH{\mathop{\mathrm{CH}}\nolimits}
\def\proetale{{pro\text{-}\acute{e}tale}}
\def\etale{{\acute{e}tale}}
\def\QCoh{\mathit{QCoh}}
\def\Ker{\mathop{\mathrm{Ker}}}
\def\Im{\mathop{\mathrm{Im}}}
\def\Coker{\mathop{\mathrm{Coker}}}
\def\Coim{\mathop{\mathrm{Coim}}}

% Boxtimes
%
\DeclareMathSymbol{\boxtimes}{\mathbin}{AMSa}{"02}

%
% Macros for moduli stacks/spaces
%
\def\QCohstack{\mathcal{QC}\!\mathit{oh}}
\def\Cohstack{\mathcal{C}\!\mathit{oh}}
\def\Spacesstack{\mathcal{S}\!\mathit{paces}}
\def\Quotfunctor{\mathrm{Quot}}
\def\Hilbfunctor{\mathrm{Hilb}}
\def\Curvesstack{\mathcal{C}\!\mathit{urves}}
\def\Polarizedstack{\mathcal{P}\!\mathit{olarized}}
\def\Complexesstack{\mathcal{C}\!\mathit{omplexes}}
% \Pic is the operator that assigns to X its picard group, usage \Pic(X)
% \Picardstack_{X/B} denotes the Picard stack of X over B
% \Picardfunctor_{X/B} denotes the Picard functor of X over B
\def\Pic{\mathop{\mathrm{Pic}}\nolimits}
\def\Picardstack{\mathcal{P}\!\mathit{ic}}
\def\Picardfunctor{\mathrm{Pic}}
\def\Deformationcategory{\mathcal{D}\!\mathit{ef}}
\usepackage{cleveref}
\crefname{originalreference}{original reference}{original references}
\Crefname{originalreference}{Original reference}{Original references}
\makeatletter
\newlabel{algebra-lemma-colimit-syntomic}{{algebra-lemma-colimit-syntomic}{source}{}{sourcepointer-1}{}}
\newlabel{algebra-lemma-colimit-syntomic@cref}{{[originalreference][0][]algebra-lemma-colimit-syntomic}{[1][0][]source}{}{}{}}
\makeatother
\crefname{originallemma}{original lemma}{original lemmas}
\Crefname{originallemma}{Original lemma}{Original lemmas}
\makeatletter
\newlabel{lemma-flat-under-gorenstein}{{lemma-flat-under-gorenstein}{source}{}{sourcepointer-2}{}}
\newlabel{lemma-flat-under-gorenstein@cref}{{[originallemma][0][]lemma-flat-under-gorenstein}{[1][0][]source}{}{}{}}
\makeatother
\crefname{originallemma}{original lemma}{original lemmas}
\Crefname{originallemma}{Original lemma}{Original lemmas}
\makeatletter
\newlabel{lemma-gorenstein-lci}{{lemma-gorenstein-lci}{source}{}{sourcepointer-3}{}}
\newlabel{lemma-gorenstein-lci@cref}{{[originallemma][0][]lemma-gorenstein-lci}{[1][0][]source}{}{}{}}
\makeatother
\crefname{originallemma}{original lemma}{original lemmas}
\Crefname{originallemma}{Original lemma}{Original lemmas}
\makeatletter
\newlabel{lemma-regular-gorenstein}{{lemma-regular-gorenstein}{source}{}{sourcepointer-4}{}}
\newlabel{lemma-regular-gorenstein@cref}{{[originallemma][0][]lemma-regular-gorenstein}{[1][0][]source}{}{}{}}
\makeatother
\crefname{originalreference}{original reference}{original references}
\Crefname{originalreference}{Original reference}{Original references}
\makeatletter
\newlabel{more-algebra-section-properties-formal-fibres}{{more-algebra-section-properties-formal-fibres}{source}{}{sourcepointer-5}{}}
\newlabel{more-algebra-section-properties-formal-fibres@cref}{{[originalreference][0][]more-algebra-section-properties-formal-fibres}{[1][0][]source}{}{}{}}
\makeatother
\title{Being Gorenstein is a property of the formal fibres}
\author{The Stacks Project Authors}
\date{stacks-project:dualizing.tex:local-snapshot-20261012; GNU FDL 1.2 or later}
\begin{document}
\maketitle
\noindent\textit{Editorial note.} All delivered mathematical text is the author's own, copied from the pinned Stacks Project chapter source sources/stacks-dualizing/pinned-original-source.tex; no statement or proof step was rewritten, reordered or removed.\par\medskip
\noindent\textit{Editorial note.} Preamble provenance: the Stacks Project keeps one single shared preamble for every chapter; this card takes it directly from the staged sources/stacks-dualizing/preamble.tex as three verbatim spans (lines 8-12, 14-18 and 154-217), so no derived file was written and nothing was added or redefined. The wrapper class and the base packages are supplied by the assembler.\par\medskip
\noindent\textit{Editorial note.} Bibliography provenance: neither the counted statement nor its proof contains a citation (this lane rescans the whole span before writing the card), so the card carries no bibliography and no entry was reconstructed.\par\medskip
\noindent\textit{Editorial note.} Selection provenance: this card was opened under the file-based control contract (CONTROL.md), with the mechanical pre-card gate run on its statement and proof interval before the card was written.\par\medskip
\noindent\textit{Allowed standard prerequisites (editorial).}\begin{itemize}
\item The permanence of the Gorenstein property under a local complete intersection extension. Invoked inside the counted proof (source lines 4427-4428). Carries the property through the relative global complete intersection used in part (A).
\item The preservation of the Gorenstein property under a flat local homomorphism with Gorenstein fibre. Invoked inside the counted proof (source lines 4443-4444 and 4447). Gives parts (C) and (D).
\item The Gorenstein property of a regular ring. Invoked inside the counted proof (source line 4442). Supplies the Gorenstein fibre in the regular case.
\end{itemize}\medskip
\ifdefined\hypertarget\hypertarget{sourcepointer-1}{}\fi
\noindent\textit{Source reference.} Original reference algebra-lemma-colimit-syntomic: sources/stacks-dualizing/pinned-original-source.tex, source lines 4426--4426. Navigation pointer only: the counted text refers to the Stacks item algebra-lemma-colimit-syntomic, which belongs to another chapter of the project and is not part of this card. The pointer carries the target's own label as its number because this chapter's staged snapshot does not contain the other chapter's tag table; its content is not reproduced here.\par\medskip
\ifdefined\hypertarget\hypertarget{sourcepointer-2}{}\fi
\noindent\textit{Source reference.} Original lemma lemma-flat-under-gorenstein: sources/stacks-dualizing/pinned-original-source.tex, source lines 4164--4164. Navigation pointer only: the counted text refers to the Stacks item lemma-flat-under-gorenstein, which is an existing labeled result of this same chapter and is not part of this card. The pointer carries the target's own label as its number and its definition line as its span; its content is not reproduced here.\par\medskip
\ifdefined\hypertarget\hypertarget{sourcepointer-3}{}\fi
\noindent\textit{Source reference.} Original lemma lemma-gorenstein-lci: sources/stacks-dualizing/pinned-original-source.tex, source lines 4136--4136. Navigation pointer only: the counted text refers to the Stacks item lemma-gorenstein-lci, which is an existing labeled result of this same chapter and is not part of this card. The pointer carries the target's own label as its number and its definition line as its span; its content is not reproduced here.\par\medskip
\ifdefined\hypertarget\hypertarget{sourcepointer-4}{}\fi
\noindent\textit{Source reference.} Original lemma lemma-regular-gorenstein: sources/stacks-dualizing/pinned-original-source.tex, source lines 3980--3980. Navigation pointer only: the counted text refers to the Stacks item lemma-regular-gorenstein, which is an existing labeled result of this same chapter and is not part of this card. The pointer carries the target's own label as its number and its definition line as its span; its content is not reproduced here.\par\medskip
\ifdefined\hypertarget\hypertarget{sourcepointer-5}{}\fi
\noindent\textit{Source reference.} Original reference more-algebra-section-properties-formal-fibres: sources/stacks-dualizing/pinned-original-source.tex, source lines 4385--4385. Navigation pointer only: the counted text refers to the Stacks item more-algebra-section-properties-formal-fibres, which belongs to another chapter of the project and is not part of this card. The pointer carries the target's own label as its number because this chapter's staged snapshot does not contain the other chapter's tag table; its content is not reproduced here.\par\medskip
\begin{lemma}
\label{lemma-formal-fibres-gorenstein}
Properties (A), (B), (C), (D), and (E) of
More on Algebra, Section \ref{more-algebra-section-properties-formal-fibres}
hold for $P(k \to R) =$``$R$ is a Gorenstein ring''.
\end{lemma}

\begin{proof}
Since we already know the result holds for Cohen-Macaulay instead
of Gorenstein, we may in each step assume the ring we have is
Cohen-Macaulay. This is not particularly helpful for the proof, but
psychologically may be useful.

\medskip\noindent
Part (A). Let $K/k$ be a finitely generated field extension.
Let $R$ be a Gorenstein $k$-algebra.
We can find a global complete intersection
$A = k[x_1, \ldots, x_n]/(f_1, \ldots, f_c)$
over $k$ such that $K$ is isomorphic to the fraction field of $A$, see
Algebra, Lemma \ref{algebra-lemma-colimit-syntomic}.
Then $R \to R \otimes_k A$ is a relative global complete intersection.
Hence $R \otimes_k A$ is Gorenstein by Lemma \ref{lemma-gorenstein-lci}.
Thus $R \otimes_k K$ is too as a localization.

\medskip\noindent
Proof of (B). This is clear because a ring is Gorenstein
if and only if all of its local rings are Gorenstein.

\medskip\noindent
Part (C). Let $A \to B \to C$ be flat maps of Noetherian rings.
Assume the fibres of $A \to B$ are Gorenstein and $B \to C$ is regular.
We have to show the fibres of $A \to C$ are Gorenstein.
Clearly, we may assume $A = k$ is a field. Then we may assume that
$B \to C$ is a regular local homomorphism of Noetherian local rings.
Then $B$ is Gorenstein and $C/\mathfrak m_B C$ is regular, in
particular Gorenstein (Lemma \ref{lemma-regular-gorenstein}).
Then $C$ is Gorenstein by
Lemma \ref{lemma-flat-under-gorenstein}.

\medskip\noindent
Part (D). This follows from Lemma \ref{lemma-flat-under-gorenstein}.
Part (E) is immediate as the condition does not refer to the ground field.
\end{proof}

\end{document}
